\documentclass[10pt]{article}

\usepackage[margin=1in]{geometry}
\usepackage{amsmath}
\usepackage{amssymb}
\usepackage{booktabs}
\usepackage{graphicx}
\usepackage{hyperref}
\usepackage[numbers,sort&compress]{natbib}

\title{Make Attention Sub-Quadratic Again}
\author{Marcel Butucea\\Independent Researcher}
\date{May 2026}

\begin{document}
\maketitle

\begin{abstract}
Full self-attention gives modern decoder-only language models their long-range
associative power, but its candidate-scoring cost grows quadratically with
sequence length. Recent systems reduce this burden either by using optimized
exact kernels, fixed sparse patterns, query-aware page selection, or
retrieval-based attention over the native key/value cache. However, vanilla
approximate-nearest-neighbor (ANN) search over a transformer's native query and
key vectors is poorly matched to the geometry of attention: queries and keys
are trained to interact through a dot product, not necessarily to occupy a
shared retrieval space.

This paper proposes a small trainable search layer for sub-quadratic attention
candidate selection. For selected frozen transformer layers, we train per-layer
query and key projections into a low-dimensional search space. These projections
are distilled from the teacher model's own attention distributions using a
combination of teacher-topK contrastive learning and distribution-level KL
distillation. At inference, dense attention over all prior keys is replaced by
retrieval of $K$ candidate keys in the learned search space, followed by
ordinary scaled-dot-product attention over only those candidates. The base model
weights are never modified.

On a 6-layer pilot using Qwen3-4B-Instruct-2507 with 4K context and WikiText-103,
a clean block-causal $d_{\mathrm{search}}=128$ run preserves full-attention
perplexity within +0.07\% at $K=128$ and +0.01\% at $K=256$. The learned search
space captures more teacher-attention mass than a Quest-style page baseline at
equal token budget, though paired perplexity is tied with Quest at $K=256$ and
slightly worse at $K=128$. A CPU FAISS/HNSW prototype tracks exact learned
retrieval on the clean evaluation slice, demonstrating ANN compatibility, but
not production wall-clock speedup. New all-layer experiments show that applying
the same substitution to all 36 layers is feasible but not yet parity: the best
all-36 checkpoint observed so far is step 500 with +3.23\% relative PPL gap, and
per-layer retrieval diagnostics identify layers 0--2 as the weakest. A
data-driven follow-up run now reserves layers 0, 1, 2, and 35 as full attention
while training layers 3--34. The central claim is therefore narrow and
falsifiable: learned search projections can make attention-relevant key
selection compatible with standard ANN retrieval, preserving model quality in a
clean pilot while replacing dense candidate scoring with a sub-quadratic
retrieval interface.
\end{abstract}

\section{Introduction}
Scaled dot-product attention computes, for a sequence of length $N$,
\begin{equation}
    A = \operatorname{softmax}\left(\frac{QK^\top}{\sqrt{d_h}} + M\right),
    \qquad O = AV,
\end{equation}
where $Q$, $K$, and $V$ are query, key, and value tensors, $d_h$ is the head
dimension, and $M$ is a causal or block-causal mask. The scoring matrix
$QK^\top$ contains $O(N^2)$ pairwise scores. Optimized kernels such as
FlashAttention~\citep{dao2022flashattention} greatly reduce memory traffic and
improve practical runtime, but the logical candidate set remains dense: each
query is scored against every eligible key.

Long-context inference exposes the weakness of this design. For each new query
token, attention over a long KV cache requires reading, scoring, and softmaxing
over a growing number of prior keys. The empirical fact that attention
distributions are often sparse suggests an alternative: identify a small set of
relevant keys, then perform exact attention only on that set. The difficult part
is not sparse attention once candidates are known; it is finding the right
candidates cheaply and robustly.

A tempting approach is to build an ANN index over native key vectors and
retrieve nearest keys for each native query vector. RetrievalAttention shows why
this fails in practice: native $Q$ and $K$ are not trained as symmetric
nearest-neighbor embeddings~\citep{retrievalattention2024}. They are trained as
two sides of a dot-product scoring system inside a transformer layer. The
resulting query/key distribution mismatch makes off-the-shelf ANN search over
native Q/K vectors unreliable unless the index itself becomes attention-aware.

This work takes the opposite route. Instead of modifying the index to compensate
for hostile native vectors, we learn a small search space in which queries and
keys are explicitly trained to be mutually retrievable. For each substituted
layer $i$, two lightweight projections map hidden states to low-dimensional
search vectors:
\begin{equation}
    q_i^s = W_{Q_i}^s h_i,\qquad k_i^s = W_{K_i}^s h_i,
\end{equation}
with $q_i^s,k_i^s\in\mathbb{R}^{d_{\mathrm{search}}}$. Retrieval is performed
in this learned space, while the final attention output still uses the frozen
model's native $Q$, $K$, and $V$ restricted to the retrieved key set.

The contribution is deliberately modest:
\begin{enumerate}
    \item A distillation objective that trains low-dimensional per-layer search
    projections to recover attention-relevant keys from a frozen teacher model.
    \item A clean block-causal evaluation showing near-parity perplexity when
    substituting learned sparse attention into 6 layers of
    Qwen3-4B-Instruct-2507.
    \item A comparison against a Quest-style page selector showing higher
    teacher-attention mass for learned retrieval, but no clean perplexity
    advantage on the current slice.
    \item A FAISS/HNSW correctness prototype showing that the learned vectors
    are compatible with an off-the-shelf ANN index.
    \item An all-36-layer feasibility experiment showing that broad substitution
    currently compounds into a +3--4\% PPL gap, plus a new data-driven reserved
    layer experiment intended to test whether full attention on the weakest edge
    layers recovers parity.
    \item A sober accounting of what is not yet shown: no measured wall-clock
    speedup, no decode-mode KV-cache integration, no long-context task
    validation, and no multi-seed confidence intervals.
\end{enumerate}

\section{Background}
\subsection{Full attention and candidate scoring}
For a single head and a single query position $t$, full causal attention
computes scores against every valid key position $j\leq t$:
\begin{equation}
    s_{tj}=\frac{q_t^\top k_j}{\sqrt{d_h}},\qquad
    \alpha_{tj}=\frac{\exp(s_{tj})}{\sum_{\ell\leq t}\exp(s_{t\ell})},
    \qquad o_t=\sum_{j\leq t}\alpha_{tj}v_j.
\end{equation}
The per-query scoring cost is $O(Nd_h)$, and the full prefill matrix costs
$O(N^2d_h)$. Reducing attention cost requires either reducing the number of
scored keys, making scoring more IO-efficient, or both.

\subsection{Sparse and retrieval-based attention}
Sparse attention methods replace the full candidate set with a smaller set
$S_t$:
\begin{equation}
    o_t=\sum_{j\in S_t}\operatorname{softmax}_{j\in S_t}
    \left(\frac{q_t^\top k_j}{\sqrt{d_h}}\right)v_j.
\end{equation}
The quality of the method depends almost entirely on whether $S_t$ contains the
keys with high teacher-attention mass. This paper uses two retrieval-quality
metrics:
\begin{equation}
    \operatorname{Recall@K}(t)=\frac{|T_K(t)\cap S_K(t)|}{\min(K,|\mathcal{C}_t|)},
\end{equation}
where $T_K(t)$ is the teacher top-$K$ set and $\mathcal{C}_t$ is the valid
causal key set, and
\begin{equation}
    \operatorname{Mass@K}(t)=\sum_{j\in S_K(t)}\alpha^{\mathrm{teacher}}_{tj}.
\end{equation}
Mass@K is the more useful metric when teacher attention is sharp, because
missing a low-probability teacher-topK key is far less important than missing a
high-probability key.

\subsection{Quest}
Quest uses query-aware page selection over the KV cache~\citep{quest2024}. Each
page stores min/max key metadata. Given a query, Quest estimates the importance
of pages and loads only the top pages for attention. It is attractive because it
is training-free, hardware-aware, and directly targets the KV-cache memory
movement bottleneck. In this paper, Quest is treated as the strongest practical
heuristic baseline currently implemented in the repository. The comparison is
not against the original optimized Quest kernel; it is against a
correctness-oriented Quest-style page selector using the same block-causal mask
and sparse gather path as learned retrieval.

\subsection{RetrievalAttention}
RetrievalAttention directly addresses long-context attention with vector
retrieval over the KV cache~\citep{retrievalattention2024}. Its key observation
is central to this work: off-the-shelf ANN indexes over native Q/K vectors are
ineffective because query and key distributions are out of distribution with
respect to each other. RetrievalAttention handles this by using attention-aware
vector search over native vectors.

This paper instead changes the vectors. The learned projections produce $q^s$
and $k^s$ in a shared retrieval space, so standard ANN machinery can be used
without making the index attention-aware. The trade-off is training:
RetrievalAttention is training-free; this method adds a small number of
trainable projection parameters.

\section{Method}
\subsection{Frozen teacher and trainable search projections}
Let $f$ be a frozen decoder-only transformer. For a selected layer $i$, let
$h_i\in\mathbb{R}^{B\times N\times d_{\mathrm{model}}}$ be the hidden state
entering the layer's self-attention module after input layer normalization. The
method attaches a SearchProjection module to that layer:
\begin{equation}
    Q_i^s = h_i W_{Q_i}^s,\qquad K_i^s = h_i W_{K_i}^s,
\end{equation}
where
\begin{equation}
    W_{Q_i}^s,W_{K_i}^s\in\mathbb{R}^{d_{\mathrm{model}}\times d_{\mathrm{search}}}.
\end{equation}
The pilot uses linear projections. For Qwen3-4B-Instruct-2507,
$d_{\mathrm{model}}=2560$. With $d_{\mathrm{search}}=128$ and six trained
layers, the search module contains approximately 3.93M trainable parameters.
The base model remains frozen.

\subsection{Teacher attention reconstruction}
Training requires the teacher's attention distribution. A naive implementation
would request \texttt{output\_attentions=True}, but this can force attention
layers onto slower eager paths. The implementation instead captures native
post-RoPE $Q$ and $K$ tensors for selected layers and reconstructs the teacher
distribution outside the model forward:
\begin{equation}
    A_i^T=\operatorname{softmax}\left(\frac{Q_iK_i^\top}{\sqrt{d_h}}+M\right).
\end{equation}
For grouped-query attention, KV heads are repeated to match query heads before
reconstruction. The teacher attention is then aggregated across heads, using
max aggregation in the current configuration:
\begin{equation}
    \bar{A}_i^T[t,j]=\max_h A_i^T[h,t,j].
\end{equation}
The implementation includes a QK reconstruction verification test that checks
agreement between reconstructed and eager attention rankings. This matters: if
the reconstructed teacher rankings are wrong, the search projections are trained
against the wrong supervision signal.

\subsection{Contrastive teacher-topK objective}
For each query token $t$, the teacher distribution selects a positive set $P_t$
consisting of the top $K_{\mathrm{pos}}$ teacher keys. Search vectors are
L2-normalized for the contrastive loss:
\begin{equation}
    \tilde{q}_t^s=\frac{q_t^s}{\|q_t^s\|_2},\qquad
    \tilde{k}_j^s=\frac{k_j^s}{\|k_j^s\|_2}.
\end{equation}
The search similarity is
\begin{equation}
    z_{tj}=\frac{(\tilde{q}_t^s)^\top \tilde{k}_j^s}{\tau}.
\end{equation}
The multi-positive InfoNCE loss is
\begin{equation}
    \mathcal{L}_{\mathrm{NCE}}(t)=
    -\log\frac{\sum_{j\in P_t}\exp(z_{tj})}
    {\sum_{j\in\mathcal{C}_t}\exp(z_{tj})},
\end{equation}
where $\mathcal{C}_t$ is the valid causal or block-causal key set. Queries with
too few valid keys and padded query positions are excluded.

\subsection{Distribution-level distillation}
The second loss aligns the full search distribution with the teacher
distribution. Search logits are computed without L2 normalization:
\begin{equation}
    r_{tj}=\frac{(q_t^s)^\top k_j^s}{\sqrt{d_{\mathrm{search}}}}.
\end{equation}
The student distribution is
\begin{equation}
    A_i^S[t,j]=\operatorname{softmax}_{j\in\mathcal{C}_t}(r_{tj}/\tau_d).
\end{equation}
The distillation loss is
\begin{equation}
    \mathcal{L}_{\mathrm{KL}}(t)=
    D_{\mathrm{KL}}(A_i^T[t,\cdot]\|A_i^S[t,\cdot])
    =\sum_{j\in\mathcal{C}_t}A_i^T[t,j]\log\frac{A_i^T[t,j]}{A_i^S[t,j]}.
\end{equation}
The total layer-averaged objective is
\begin{equation}
    \mathcal{L}=\alpha\frac{1}{|\mathcal{I}|}\sum_{i\in\mathcal{I}}\mathcal{L}_{\mathrm{NCE}}^i
    +\beta\frac{1}{|\mathcal{I}|}\sum_{i\in\mathcal{I}}\mathcal{L}_{\mathrm{KL}}^i,
\end{equation}
with $\alpha=\beta=1$ in the pilot.

\subsection{Inference-time substitution}
At inference, selected attention layers are monkey-patched. For a substituted
layer:
\begin{enumerate}
    \item Compute the native $Q$, $K$, and $V$ exactly as the base model would,
    including q\_norm/k\_norm and RoPE.
    \item Compute $Q^s$ and $K^s$ from the same hidden states.
    \item Retrieve $K_{\mathrm{retrieve}}$ candidate keys per query in the
    learned search space.
    \item Gather native $K$ and $V$ at those retrieved positions.
    \item Compute normal scaled-dot-product attention over the gathered
    candidate set only.
    \item Apply the original output projection.
\end{enumerate}
The final sparse attention is therefore not an approximation in value space.
The approximation is in candidate selection only. For a query $t$ and retrieved
set $S_t$, the substituted attention output is
\begin{equation}
    \hat{o}_t=\sum_{j\in S_t}
    \frac{\exp(q_t^\top k_j/\sqrt{d_h})}{\sum_{\ell\in S_t}\exp(q_t^\top k_\ell/\sqrt{d_h})}v_j.
\end{equation}

\subsection{Block-causal packing}
The initial packed experiments were leakage-confounded: tokens from different
packed documents could attend across document boundaries. The clean evaluation
uses segment-level block-causal masking. Each packed document receives a
segment id, position ids are reset at segment boundaries, and the attention mask
allows a query token to attend only to earlier tokens in the same segment:
\begin{equation}
    M_{tj}=0\quad\mathrm{iff}\quad
    \operatorname{segment}(t)=\operatorname{segment}(j)\ \mathrm{and}\ j\leq t,
\end{equation}
and $M_{tj}=-\infty$ otherwise. Retrieval, loss masking, Recall@K, Mass@K, and
sparse attention all use this same eligibility relation.

\section{Complexity analysis}
This section analyzes candidate scoring, not the full attention layer. Once
$K$ candidates are selected, sparse attention still costs $O(Kd_h)$ per query
for native scoring and value aggregation.

For full attention, the per-query candidate-scoring proxy is
\begin{equation}
    C_{\mathrm{full}}(N)=Nd_h.
\end{equation}
With $d_h=128$,
\begin{equation}
    C_{\mathrm{full}}(N)=128N.
\end{equation}
For Quest-style page selection with page size $P=16$, using min/max metadata
over native keys,
\begin{equation}
    C_{\mathrm{Quest}}(N)=\frac{N}{P}\cdot 2d_h=16N.
\end{equation}
For learned HNSW retrieval, using HNSW parameters $M=32$,
$ef_{\mathrm{search}}=64$, and $d_{\mathrm{search}}=128$, the deterministic
operation-count proxy used in the repository is
\begin{equation}
    C_{\mathrm{HNSW}}(N)=M\cdot ef_{\mathrm{search}}\cdot \log_2(N)\cdot d_{\mathrm{search}}.
\end{equation}
This gives
\begin{equation}
    C_{\mathrm{HNSW}}(N)=32\cdot64\cdot128\cdot\log_2(N).
\end{equation}
This is not a measured GPU runtime model and should not be read as a guarantee
for HNSW. It is a candidate-scoring proxy. Under these constants, Quest is
cheaper below the few-hundred-thousand-token regime, while learned HNSW becomes
cheaper at very long contexts. The repository's deterministic scaling table
places the rough crossover against Quest around the 300K-token range and gives
an approximately 3x scoring-proxy advantage for learned HNSW at 1M context.

The practical implication is sharp: this method is not automatically a faster
replacement for Quest at 4K or 32K context. The argument becomes interesting
when context length is very large and candidate selection must avoid linear
scans over pages or tokens.

\section{Experimental setup}
\subsection{Model}
The pilot uses Qwen/Qwen3-4B-Instruct-2507, a 36-layer decoder-only model. The
configuration used by the repository records hidden\_size=2560, 32 query heads,
8 KV heads, grouped-query attention, head\_dim=128, and a native context length
of 262K. The pilot substitutes six layers:
\begin{equation}
    \mathcal{I}=\{4,8,12,16,20,24\}.
\end{equation}

\subsection{Data and training}
The training data is WikiText-103 raw text. The pilot sequence length is 4096,
batch size is 8, and the main 6-layer runs train for 1000 steps. The
recommended clean configuration is the packed block-causal
$d_{\mathrm{search}}=128$ run. Search projections are trained while all
base-model weights remain frozen.

\subsection{Metrics}
The paper reports:
\begin{itemize}
    \item Perplexity under full attention and substituted sparse attention.
    \item Relative PPL gap:
    $(\mathrm{PPL}_{\mathrm{sparse}}-\mathrm{PPL}_{\mathrm{full}})/
    \mathrm{PPL}_{\mathrm{full}}$.
    \item Recall@K against teacher topK.
    \item Mass@K, the teacher probability mass captured by the retrieved set.
    \item FAISS filler rate for the prototype ANN path.
    \item Paired NLL deltas for learned-vs-Quest comparison where available.
\end{itemize}
PPL is the primary model-quality metric. Mass@K is the primary
retrieval-fidelity metric. Recall@K is reported but is less reliable when
teacher attention has a long low-probability tail.

\section{Results}
\subsection{Capacity ablation on packed WikiText-103}
The packed $d_{\mathrm{search}}$ ablation should be interpreted as a capacity
comparison, not as the clean headline result, because the original packed path
had cross-document leakage. At $K=128$:
\begin{table}[h]
\centering
\begin{tabular}{rrrrrr}
\toprule
$d_{\mathrm{search}}$ & Trainable params & Learned mass@K & Raw-QK oracle mass & Learned/oracle & Final PPL gap\\
\midrule
64 & 1.97M & 0.492 & 0.488 & 1.01x & +2.39\%\\
128 & 3.93M & 0.503 & 0.488 & 1.03x & -1.81\%\\
256 & 7.86M & 0.509 & 0.488 & 1.04x & -1.85\%\\
\bottomrule
\end{tabular}
\end{table}
The capacity trend is monotonic but saturating. $d_{\mathrm{search}}=128$
captures almost all of the $d_{\mathrm{search}}=256$ retrieval quality at half
the trainable parameter count. For that reason, $d_{\mathrm{search}}=128$ is
the recommended pilot capacity.

\subsection{Clean block-causal \texorpdfstring{$d_{\mathrm{search}}=128$}{dsearch=128}}
The clean block-causal run removes cross-document leakage. On the 16-batch clean
evaluation slice, full-attention PPL is 30.44.
\begin{table}[h]
\centering
\begin{tabular}{rrrrr}
\toprule
$K$ & Recall@K & Mass@K & Sparse PPL & Relative PPL gap\\
\midrule
128 & 0.744 & 0.787 & 30.47 & +0.07\%\\
256 & 0.879 & 0.953 & 30.45 & +0.01\%\\
512 & n/a & n/a & 30.45 & +0.01\%\\
\bottomrule
\end{tabular}
\end{table}
K=512 has no meaningful average recall/mass on this WikiText slice because
almost no same-segment queries have 512 valid causal keys. The PPL result is
still useful: filler slots are masked out of the sparse-attention softmax, and
the substituted path remains effectively at full-attention parity.

Clean block-causal per-layer Mass@K at $K=128$:
\begin{table}[h]
\centering
\begin{tabular}{rrr}
\toprule
Layer & Raw-QK oracle mass & Learned d128 mass\\
\midrule
4 & 0.956 & 0.950\\
8 & 0.977 & 0.976\\
12 & 0.970 & 0.977\\
16 & 0.964 & 0.970\\
20 & 0.970 & 0.983\\
24 & 0.978 & 0.984\\
Average & 0.969 & 0.973\\
\bottomrule
\end{tabular}
\end{table}
The clean result changes the earlier interpretation. Under segment-isolated
masking, early trained layers are not uniquely hard; all six tested layers have
high teacher-attention mass, and the learned projection matches or slightly
exceeds raw-QK retrieval across the tested set.

\subsection{Quest-style page baseline}
A Quest-style min/max page baseline was evaluated with page size 16, native
post-RoPE Q/K, the same block-causal token eligibility mask, and the same
sparse-attention gather path.
\begin{table}[h]
\centering
\begin{tabular}{lrrrrr}
\toprule
Method & $K$ & Recall@K & Mass@K & PPL & Relative PPL gap\\
\midrule
Learned search exact & 128 & 0.744 & 0.787 & 30.47 & +0.07\%\\
Quest-style page & 128 & 0.669 & 0.727 & 30.41 & -0.11\%\\
Learned search exact & 256 & 0.879 & 0.953 & 30.45 & +0.01\%\\
Quest-style page & 256 & 0.838 & 0.909 & 30.45 & +0.03\%\\
\bottomrule
\end{tabular}
\end{table}
Learned search recovers more teacher-attention mass at equal $K$, especially at
$K=128$. However, perplexity does not currently show a clean advantage over
Quest. A paired 32-batch NLL evaluation gives the sharper reading:
\begin{table}[h]
\centering
\begin{tabular}{rrrrl}
\toprule
$K$ & Full PPL & Learned PPL & Quest PPL & Learned--Quest NLL delta, 95\% bootstrap CI\\
\midrule
128 & 28.03 & 28.07 & 28.01 & +0.00205 $[+0.00160,+0.00251]$; Quest slightly better\\
256 & 28.03 & 28.04 & 28.04 & -0.00005 $[-0.00029,+0.00018]$; statistical tie\\
\bottomrule
\end{tabular}
\end{table}
The honest claim is therefore not ``learned search beats Quest in PPL.'' The
current claim is ``learned search provides higher retrieval fidelity and ANN
compatibility; PPL is tied with Quest at $K=256$ and slightly worse at $K=128$
on the paired WikiText slice.''

\subsection{FAISS/HNSW compatibility}
The first FAISS prototype used a global index followed by segment filtering,
which caused pathological filler rates. The corrected clean path builds
per-segment indexes when a block-causal mask is present. With that fix, CPU
FAISS/HNSW tracks exact learned search on the clean 16-batch evaluation slice:
\begin{table}[h]
\centering
\begin{tabular}{lrrr}
\toprule
Method & $K$ & PPL & Relative PPL gap\\
\midrule
Learned exact & 128 & 30.47 & +0.07\%\\
Learned FAISS/HNSW & 128 & 30.47 & +0.09\%\\
Learned exact & 256 & 30.45 & +0.01\%\\
Learned FAISS/HNSW & 256 & 30.46 & +0.04\%\\
\bottomrule
\end{tabular}
\end{table}
The filler rate is expected for short same-segment prefixes where fewer than
$K$ valid causal keys exist. Filler slots are masked out of the
sparse-attention softmax. These results demonstrate compatibility with
off-the-shelf ANN retrieval, not production efficiency.

\subsection{Dynamic-index proxy}
The repository currently implements a prefill-only path with
\texttt{use\_cache=False}. A true decode-time dynamic-index benchmark still
requires KV-cache integration. As a proxy, the dynamic-index experiment splits
clean block-causal sequences into a prefill prefix and a decode-like suffix. It
compares:
\begin{itemize}
    \item dynamic retrieval, where suffix queries can retrieve from all
    same-segment prior keys;
    \item static retrieval, where suffix queries can retrieve from prefill keys
    plus a 256-token local suffix window but not older suffix keys.
\end{itemize}
At $K=128$, prefill length 1024, local window 256, and 8 eval batches:
\begin{table}[h]
\centering
\begin{tabular}{lr}
\toprule
Setting & Teacher mass captured\\
\midrule
Dynamic proxy & 0.972\\
Static proxy & 0.928\\
Static teacher mass available & 0.954\\
Dynamic -- static & +0.044\\
\bottomrule
\end{tabular}
\end{table}
This does not establish decode accuracy or latency. It only shows that a frozen
prefill-plus-local index loses measurable teacher-attention mass relative to a
dynamic index on decode-like suffix queries.

\subsection{All-layer and reserved-layer substitution}
After the 6-layer clean pilot, we launched a stronger full-substitution
experiment with $d_{\mathrm{search}}=128$ on all 36 attention layers under the
same packed block-causal methodology. This run uses batch size 2 with gradient
accumulation 4, checkpointing every 25 steps, and the same $K=128$ training eval
slice. The run reached step 800 before we stopped it to analyze checkpoints and
start a data-driven reserved-layer follow-up.

\begin{table}[h]
\centering
\caption{All-36 $d_{\mathrm{search}}=128$ training trajectory.}
\begin{tabular}{rrr}
\toprule
Step & Recall@K eval & Relative PPL gap\\
\midrule
250 & 0.805 & +6.271\%\\
500 & 0.816 & +3.227\%\\
750 & 0.817 & +3.964\%\\
\bottomrule
\end{tabular}
\end{table}

The all-36 result is good but not a parity result. Broad substitution preserves
substantial retrieval fidelity, but the per-layer approximation errors compound
into a +3--4\% PPL gap. Step 500 is the best checkpoint observed by PPL. Steps
500, 750, and 800 have essentially identical Mass@K, so the non-monotonic PPL
trajectory is not explained by a large retrieval-mass change.

\begin{table}[h]
\centering
\caption{All-36 step-500 per-layer Mass@K at $K=128$. Delta is learned minus raw-QK.}
\begin{tabular}{rrrrrrrr}
\toprule
Layer & Raw & Learned & Delta & Layer & Raw & Learned & Delta\\
\midrule
0 & 0.922 & 0.780 & -0.142 & 18 & 0.965 & 0.972 & +0.007\\
1 & 0.918 & 0.851 & -0.067 & 19 & 0.961 & 0.968 & +0.007\\
2 & 0.939 & 0.899 & -0.040 & 20 & 0.959 & 0.975 & +0.016\\
3 & 0.939 & 0.924 & -0.015 & 21 & 0.966 & 0.979 & +0.014\\
4 & 0.944 & 0.933 & -0.011 & 22 & 0.963 & 0.970 & +0.007\\
5 & 0.964 & 0.947 & -0.017 & 23 & 0.979 & 0.984 & +0.005\\
6 & 0.956 & 0.936 & -0.020 & 24 & 0.971 & 0.978 & +0.007\\
7 & 0.982 & 0.982 & +0.000 & 25 & 0.986 & 0.988 & +0.002\\
8 & 0.971 & 0.970 & -0.001 & 26 & 0.978 & 0.985 & +0.008\\
9 & 0.959 & 0.976 & +0.017 & 27 & 0.978 & 0.983 & +0.005\\
10 & 0.974 & 0.970 & -0.004 & 28 & 0.979 & 0.985 & +0.005\\
11 & 0.976 & 0.975 & -0.001 & 29 & 0.982 & 0.987 & +0.005\\
12 & 0.961 & 0.969 & +0.008 & 30 & 0.988 & 0.986 & -0.002\\
13 & 0.971 & 0.971 & +0.000 & 31 & 0.984 & 0.984 & -0.001\\
14 & 0.973 & 0.973 & -0.000 & 32 & 0.979 & 0.979 & +0.000\\
15 & 0.968 & 0.972 & +0.004 & 33 & 0.977 & 0.970 & -0.007\\
16 & 0.956 & 0.962 & +0.006 & 34 & 0.976 & 0.960 & -0.016\\
17 & 0.959 & 0.966 & +0.007 & 35 & 0.980 & 0.967 & -0.013\\
\midrule
\multicolumn{5}{r}{Average} & 0.966 & 0.960 & -0.006\\
\bottomrule
\end{tabular}
\end{table}

The diagnostic is clear: the weak layers are concentrated at the beginning of
the network. Layer 0 is the main outlier, layer 1 is weak, and layer 2 is
borderline. Layer 35 is mildly below raw-QK but not the primary failure mode.
Based on this, the next experiment reserves layers 0, 1, 2, and 35 as full
attention and trains learned retrieval for layers 3--34. That run is named
\texttt{all32-d128-block-causal-reserve-0-1-2-35}. At the time of this draft,
the run is active; the first checkpoint at step 25 has been saved. This is the
right next test because it asks whether broad substitution fails because of a
few edge layers or because small per-layer errors compound even when every
individual layer has high Mass@K.

\section{Discussion}
\subsection{What the result proves}
The clean result proves that learned low-dimensional search projections can
preserve the quality of selected transformer layers when sparse attention
candidates are selected in the learned space. It also shows that standard HNSW
retrieval can approximate exact learned retrieval closely enough on the clean
slice.

This is important because it separates two problems:
\begin{enumerate}
    \item Can the model's attention-relevant keys be represented in a small
    shared retrieval space?
    \item Can a production kernel retrieve from that space faster than dense
    attention or page selection?
\end{enumerate}
The pilot gives evidence for the first question. It does not yet solve the
second.

\subsection{What the result does not prove}
The current repository does not prove measured wall-clock speedup. The FAISS
path builds CPU indexes per forward pass, performs GPU-to-CPU transfers, and
uses dense-style tensor expansion internally for gather. This is a correctness
prototype.

The current repository also does not prove full-model parity. The clean 6-layer
result is near parity, but the all-36 run observed so far has a +3--4\% PPL
gap. The new all32 reserved-layer run is designed to test whether this is mostly
an edge-layer issue.

The current repository does not prove long-context task quality. WikiText
perplexity is a useful first sanity check but is not a substitute for LongBench,
RULER, passkey retrieval, needle-in-haystack, or downstream generation
benchmarks.

Finally, the clean comparison does not beat Quest in perplexity. Learned search
has higher Mass@K, but Quest is slightly better at $K=128$ and tied at $K=256$
on the paired NLL evaluation.

\subsection{Why learned retrieval may still matter}
Quest is a strong short- and medium-context baseline because page metadata gives
a cheap query-aware score. Under the repository's scoring proxy, Quest is
cheaper than learned HNSW below a few hundred thousand tokens. Learned ANN
retrieval becomes interesting when:
\begin{itemize}
    \item context length reaches hundreds of thousands to millions of tokens;
    \item the retrieval index can live on GPU or in a low-latency paged memory
    system;
    \item the index can be incrementally updated during decode;
    \item higher Mass@K translates into task quality at lower $K$ or under
    harder long-context distributions.
\end{itemize}
The next experiments should therefore not merely repeat WikiText. They should
test whether learned retrieval allows lower $K$ than Quest on tasks with sparse,
long-range dependencies.

\section{Related work}
\subsection{Property-level comparison}
The most useful comparison is not only to Quest. Quest is the closest practical
page-selection baseline, but Reformer is closer in architectural shape:
candidate selection is performed by an ANN-like mechanism and attention is then
computed over a restricted set. Performer is also genuinely sub-quadratic, but
it changes the attention computation itself by approximating the softmax kernel
with random features. Longformer and BigBird use sparse patterns that are cheap
but mostly fixed rather than query-adaptive. These distinctions matter because a
long-context method can fail in different ways: it may be sub-quadratic but
untrained, trained but still asymptotically dense, query-aware but linear in the
number of pages, or fast but unable to represent sharp retrieval.

\begin{table}[h]
\centering
\caption{Qualitative comparison of sparse and efficient attention mechanisms.}
\label{tab:related-properties}
\begin{tabular}{llllll}
\toprule
Method & Selection mechanism & Query-aware & Trained & Asymptotic & Exact softmax\\
\midrule
Full attention & all keys & n/a & n/a & $O(N^2)$ & yes\\
Reformer & LSH hashing & yes & no & $O(N\log N)$ & over bucket\\
Performer & random features & n/a & no & $O(N)$ & no\\
BigBird & window + random + global & mostly no & no & $O(N)$ & over pattern\\
Longformer & sliding window + global & mostly no & no & $O(N)$ & over pattern\\
NSA-style methods & block compression/selection & partial & partial & $O(N^2)$ proxy & yes\\
Quest & min/max page heuristic & yes & no & $O(N)$ & over pages\\
This work & trained low-dim retrieval & yes & yes & $O(N\log N)$ & over retrieved set\\
\bottomrule
\end{tabular}
\end{table}

This table is a positioning argument, not a completed empirical victory. The
current experiments establish the learned-retrieval row for a clean six-layer
pilot and show that all-layer substitution is feasible but not yet parity. The
active all32 reserved-layer run is meant to test whether a near-whole-model
version can preserve the six-layer quality result. Until that run and
long-context tasks are complete, the table should be read as the design target
and not as proof of production superiority.

\paragraph{Reformer as the closest asymptotic ancestor.}
Reformer's LSH attention can be viewed as untrained ANN-style candidate
selection: hash the sequence, attend within nearby buckets, and avoid scoring
every key~\citep{kitaev2020reformer}. The appealing property is genuine
sub-quadratic scaling. The weakness is that random hash functions do not know
which keys the model's attention actually needs. Native Q/K vectors are also not
guaranteed to be good nearest-neighbor embeddings. This work keeps Reformer's
``retrieve then attend'' shape but replaces untrained hashing with projections
distilled from the teacher model's attention distribution.

\paragraph{Performer and smooth attention approximations.}
Performer approximates softmax attention with positive orthogonal random
features, giving linear-time attention in principle. This is a different
failure mode from sparse retrieval. Linear kernel approximations replace the
softmax with a smoother low-rank computation, so they can lose the sharp
selectivity needed for precise key-value retrieval. Our method does not
approximate the softmax inside the selected set; it preserves native scaled
dot-product attention over retrieved keys. The approximation is only in
candidate selection.

\paragraph{Fixed sparse patterns.}
Longformer and BigBird reduce cost with local windows, global tokens, random
connections, or structured sparse patterns~\citep{beltagy2020longformer,zaheer2020bigbird}.
These patterns are efficient and useful, but they are not fully query-aware in
the sense used here: the selected remote keys are not chosen because a specific
query token currently needs them. Query-aware methods such as Quest and learned
retrieval are better matched to long-context retrieval tasks where the relevant
token may be far away and content-dependent.

\paragraph{NSA-style block selection.}
Recent native sparse attention systems based on block compression and block
selection can deliver strong constant-factor speedups, but if block relevance is
computed by comparing queries to compressed block representations across the
sequence, the selection branch can remain asymptotically quadratic in sequence
length up to a block-size constant. That does not make such methods unimportant:
constant factors dominate at many practical lengths. It does mean they occupy a
different asymptotic regime from ANN retrieval, where the intended candidate
selection cost is sub-quadratic.

\paragraph{Transformers and exact attention.}
The Transformer introduced scaled dot-product attention as the central
mechanism for sequence modeling~\citep{vaswani2017attention}. FlashAttention
later showed that exact attention can be made substantially more IO-efficient
using tiling and SRAM-aware computation~\citep{dao2022flashattention}. This
paper is complementary: it does not optimize exact dense attention; it changes
the candidate-selection problem.

\paragraph{Fixed sparse attention.}
Longformer, BigBird, Routing Transformer, Reformer, and related systems reduce
attention cost using fixed, random, clustered, local, or hashing-based sparse
patterns~\citep{beltagy2020longformer,zaheer2020bigbird,kitaev2020reformer,roy2021routing}.
These approaches avoid dense scoring but may miss query-specific important keys
unless the pattern is well matched to the task. Learned search projections are
query-specific and layer-specific.

\paragraph{Query-aware KV selection.}
Quest selects KV pages using query-aware min/max metadata and is a strong
practical baseline~\citep{quest2024}. The present method differs by learning a
retrieval embedding rather than relying on native Q/K page bounds. The current
evidence favors Quest on practical PPL at $K=128$ and ties at $K=256$, while
learned retrieval captures more teacher mass and offers an ANN-compatible route
for very long contexts.

\paragraph{RetrievalAttention.}
RetrievalAttention is closest in motivation~\citep{retrievalattention2024}. It
observes that native Q/K vectors are poorly suited for vanilla ANN due to
query/key distribution mismatch, and it fixes this with attention-aware search.
This paper fixes the mismatch by training a small shared search space. In one
sentence: RetrievalAttention adapts the index to native Q/K; this work adapts
Q/K-like search vectors to the index.

\paragraph{ANN indexing.}
HNSW and FAISS provide practical approximate nearest-neighbor search
machinery~\citep{malkov2018hnsw,johnson2019billion}. The current implementation
uses FAISS HNSW as a correctness check. A deployable system would need a
GPU-resident or cache-integrated index, not per-forward CPU index construction.

\section{Limitations and next experiments}
The following items are required before this work should be presented as more
than a promising research prototype:
\begin{enumerate}
    \item Multi-seed confidence intervals for the clean block-causal result.
    \item Larger eval slices and stricter paired testing against Quest.
    \item A four-way decomposition: full attention vs exact native-QK topK vs
    exact learned-search topK vs FAISS learned-search topK.
    \item Long-context task evaluation on LongBench, RULER, passkey retrieval,
    and needle-in-haystack.
    \item Completion and evaluation of the active all32 reserved-layer run.
    \item Decode-mode KV-cache integration with incremental index updates.
    \item GPU-resident retrieval and fused sparse gather-attention kernels.
    \item Measured wall-clock latency and memory footprint against
    FlashAttention/SDPA and optimized Quest.
    \item Ablation over $d_{\mathrm{search}}$, $K_{\mathrm{pos}}$,
    KL-vs-contrastive weighting, head aggregation, and per-head vs shared-head
    retrieval.
    \item Cross-model validation beyond Qwen3-4B-Instruct-2507.
\end{enumerate}

\section{Conclusion}
This paper proposes learning a small per-layer search space for attention
candidate retrieval. The method keeps the base language model frozen, trains
lightweight query/key search projections from teacher attention, retrieves
candidate keys in the learned space, and performs normal attention over the
retrieved native KV vectors.

The clean pilot result is encouraging but narrow: six substituted layers of
Qwen3-4B-Instruct-2507 preserve WikiText perplexity under block-causal masking,
and FAISS/HNSW closely tracks exact learned retrieval. Compared with Quest,
learned retrieval captures more teacher-attention mass but does not yet deliver
a clean perplexity win. The first all-36 result shows that broad substitution is
feasible but still costs +3--4\% relative PPL under the current hyperparameters.
The result should therefore be framed as an ANN-compatibility and
retrieval-fidelity contribution, not as a completed fast-attention system.

The path to a stronger claim is clear: finish the all32 reserved-layer run, run
long-context tasks, integrate decode-mode dynamic indexing, and demonstrate
measured latency gains with a GPU-resident retrieval kernel. Until then, the
correct title is aspirational but defensible: Make Attention Sub-Quadratic
Again.

\bibliographystyle{plainnat}
\bibliography{references}

\end{document}
